\section{Examples and Definitions}
\label{sec:examples-and-definitions}

There are two ways to view groups, \emph{algebra} and \emph{symmetry}. We first build intuition by looking at an example in symmetry.

\begin{example}[Equilateral Triangle]
    An equilateral triangle has rotational symmetries and reflective symmetries. Note that there is the \emph{trivial symmetry} which is the \emph{identity}, that does not move the triangle and keeps it in place.

    An equilateral triangle has 6 symmetries.
\end{example}

Natural questions arise: what are the symmetries of a square, a regular pentagon, or in general, a regular \(n\)-gon?

We can define 'multiplication' on symmetries, specifically by composing the symmetries:

We note the following features (and non-features) arise in the previous example. \bigbreak

\noindent
\textbf{Important features.}
\begin{itemize}
    \item There is an \emph{identity}.
    \item Every symmetry has an \emph{inverse}.
    \item Symmetries can be \emph{composed} (to form another symmetry).
    \item Composition of symmetries are \emph{associative}.
\end{itemize}

\noindent
\textbf{Non-feature.} The symmetries need not commute.

\subsection{Introduction to the Algebraic Standpoint}
\label{sec:algebraic-standpoint}

\begin{definition}[Binary Operation]
    A \emph{binary operation} on a set \(X\) is a function
    \[
        \cdot : X \times X \to X
    \]
    where \(X \times X\) is the \emph{Cartesian product} given by
    \[
        X \times X = \set{\circPara{x_1, x_2}}{x_1, x_2 \in X}.
    \]
    
    We denote
    \[
        \circPara{x_1, x_2} \mapsto x_1 \cdot x_2.
    \]
\end{definition}

\begin{definition}[Group]
    A \emph{group} is a triple
    \[
        \circPara{G, \cdot, e}
    \]
    where
    \begin{itemize}
        \item \(G\) is a set,
        \item \(\cdot\) is a binary operation on \(G\), and
        \item \(e\) is an element of \(G\), denoted as \(e \in G\),
    \end{itemize}
    such that the following 4 axioms are satisfied:
    \begin{enumerate}
        \item \emph{Closure}. For all \(a, b \in G\), \(a \cdot b \in G\).
        \item \emph{Associativity}. For all \(a, b, c \in G\), \(
        \circPara{\circPara{a \cdot b} \cdot c} = \circPara{a \cdot \circPara{b \cdot c}}\).
        \item \emph{(Right) Identity}. For all \(a \in G\), \(a \cdot e = a\).
        \item \emph{(Right) Inverses}. For all \(a \in G\), there exists some \(b \in G\), such that \(a \cdot b = e\).
    \end{enumerate}
\end{definition}

\begin{remark}[Remark on Closure]
    The closure axiom is redundant, since \(\cdot\) being a binary operation means that \(G\) has to be closed under \(\cdot\)
\end{remark}

A group is an algebraic structure with one operation.

\begin{example}[Equilateral Triangle]
    The symmetries of an equilateral triangle forms a group, with respect to composition, and the identity symmetry is the identity element.
\end{example}

\begin{example}[Integers]
    \(\circPara{\ZZ, +, 0}\) forms a group.
\end{example}

\begin{proposition}[Properties on Inverses and Identities]
    Let \(\circPara{G, \cdot, e}\) be a group. Let \(a, b, b', e' \in G\).

    \begin{enumerate}
        \item \emph{Right-inverses are also left-inverses.} If \(a \cdot b = e\), then \(b \cdot a = e\).
        \item \emph{A right-identity is also a left-identity.} \(e \cdot a = a\).
        \item \emph{Inverses are unique.} If \(a, b, b'\) are such that \(a \cdot b = e\) and \(a \cdot b' = e\), then \(b = b'\).
        \item \emph{The identity is unique.} If \(a \cdot e' = a\), then \(e' = e\).
    \end{enumerate}
\end{proposition}